The fixed syntactic order is not needed for the semantic execution
interface. Let a \emph{well-founded dependent code} be an inductively
generated object
\[
 G=\End\langle R_i\rangle_i
 \quad\text{or}\quad G=\Q_r(a:A).H(a),
\]
where $H$ is a family of codes, not merely a syntactic continuation
whose binder positions are fixed. Its next owner, answer type, number
of moves, and terminal return types may all depend on $a$. The inductive
formation condition requires well-foundedness; it does not require a
uniform finite depth.

At a terminal, both the local interface $L_i$ and the full-strategy
interface $F_i$ are $R_i$. At a binder, define
\[
\begin{array}{c|cc}
 &i=r&i\ne r\\\hline
 L_i(G)&\displaystyle\sum_{a:A}L_i(H(a))
       &\displaystyle\prod_{a:A}L_i(H(a))\\[3pt]
 F_i(G)&\displaystyle A\times\prod_{a:A}F_i(H(a))
       &\displaystyle\prod_{a:A}F_i(H(a)).
\end{array}
\]
The difference is the owner's retention of every alternative in $F_i$.
The trace, return family, and complete execution record have exactly
the dependent definitions of \cref{sec:interaction}.

\begin{theorem}[Value-dependent execution]\label{dep:thm:dynamic-execution}
Every well-founded dependent code has well-typed endpoint and full-policy
execution functions. Reducing a full profile to its selected endpoint plans
preserves the entire shared trace and all dependent terminal returns:
\[
 \Run_G^{\rm full}(\vec\sigma)
   =\Run_G(\operatorname{reduce}_G(\vec\sigma)).
\]
Selecting a legal answer always descends in a well-founded child
relation. There are such codes with complete executable profiles of
arbitrarily large finite trace length and no uniform finite depth bound.
\end{theorem}
\begin{proof}
Define reduction by well-founded recursion. At an owned binder keep
the owner's selected $a$ and reduce its $a$-continuation; at every
other endpoint reduce the entire receiving family. Both executions
select the same $a$, by the owned component of the profile, and pass
the same $a$ to every endpoint. The induction hypothesis equates their
complete child records, including the return types indexed by the
child trace. Adding the common first answer gives equality of the
dependent execution records. At a terminal both return the same
record directly. The child relation is well founded by induction on
the code.

For unbounded length, Alice first chooses $n\in\mathbb N$. The
continuation has $n$ moves alternating between Bob and Carol, starting
with Bob when $n$ is even and Carol when $n$ is odd, with
an action type $\operatorname{Fin}(j+1)$ when $j+1$ moves remain.
Each menu is inhabited; choosing zero at every subsequent legal
history gives a complete profile with $n+1$ decisions. A uniform
depth bound would bound all these lengths, a contradiction.
\end{proof}

\paragraph{Private boundaries and legal views.}
For private control, the answer family must be indexed by the acting
player's legal view, not by a hidden history or the address of a
code branch. Otherwise merely presenting the menu may reveal hidden
data. This requirement also determines how full policies compose.

A view arena is a well-founded tree with constructors $\operatorname{done}(b)$
and $\operatorname{ask}(v,H)$, where $H:A(v)\to\mathsf{Arena}$ and the
answer type $A(v)$ depends only on the decision view $v$. The same
view may occur at several histories. The result type can be a dependent
record packaged with its boundary context. A complete policy on an
arena is a section of $A$ on its set of \emph{all legal} decision views,
including those inconsistent with its own earlier choices. Giving each
view an owner makes a whole-player deviation a change of this section
only at that player's views.

Let $G$ be a prefix arena and $H_b$ its continuation at result $b$.
Write $P$ for the prefix's legal views and
$Q=\bigcup_{b\text{ a legal exit of }G}\operatorname{Views}(H_b)$.
For a domain $D$ put
$\operatorname{Pol}(D)=\prod_{v\in D}A(v)$.

\begin{theorem}[Dynamic boundary gluing]\label{dep:thm:dynamic-gluing}
The graft $G;H$ has legal view set $P\cup Q$, and its full policy
space is canonically
\[
 \operatorname{Pol}(P\cup Q)\cong
 \operatorname{Pol}(P)\mathbin{\times}_{\operatorname{Pol}(P\cap Q)}
 \operatorname{Pol}(Q).
\]
Restriction and gluing preserve execution: run the restricted prefix
policy first, then the same full suffix policy at the exit reached.
Whole-player deviations correspond exactly to compatible pairs of
component deviations. Grafting is associative and terminal return
is its identity. If $P\cap Q=\varnothing$, the policy space is the
ordinary product. In general that product is incorrect.
\end{theorem}
\begin{proof}
Induction on $G$ shows that every legal request in its graft lies
either in $P$ or in a continuation at a legal exit, and that all such
requests occur. A section on $P\cup Q$ restricts to sections agreeing
on the overlap. Conversely two agreeing sections glue by using the
first on $P$ and the second on $Q$; agreement makes this well defined.
The operations are inverse. The execution law follows by induction:
the root request has the same answer under restriction, and the
induction hypothesis handles the selected child. Equality of policies
at every other player's view restricts to the two domains, and the
converse follows since they cover the graft's legal views. Induction
also proves associativity of grafting and the return identity.
Disjoint domains impose no compatibility condition. Without that
condition a shared Boolean view could be assigned false by the
prefix and true by the suffix, which no single complete policy can do.
\end{proof}

This gives a semantic account of value-dependent move orders and
private boundary composition. Interpreting a syntactic disclosure
calculus into it still requires audience-derived views, dependent
context projections, and the insertion/substitution and sealing
properties of that calculus. Arbitrary view labels do not establish
a hiding theorem. The construction composes full policies and
executions; it selects no equilibrium at an intermediate boundary.
